% BHU PYQ -- Manually transcribed & error-corrected
\documentclass[11pt,a4paper]{article}

\usepackage[a4paper,top=2.5cm,bottom=2.5cm,left=2.8cm,right=2.8cm,headheight=14pt]{geometry}
\usepackage{amsmath,amssymb}
\usepackage[T1]{fontenc}
\usepackage[utf8]{inputenc}
\usepackage{lmodern}
\usepackage{microtype}
\usepackage{fancyhdr}
\usepackage{enumitem}
\usepackage{hyperref}
\usepackage{lastpage}
\usepackage{parskip}

\hypersetup{colorlinks=true,urlcolor=black,linkcolor=black,
  pdftitle={BPT-602: Solid State Physics -- BHU PYQ},pdfauthor={Banaras Hindu University}}

\pagestyle{fancy}\fancyhf{}
\renewcommand{\headrulewidth}{0.4pt}\renewcommand{\footrulewidth}{0.4pt}
\fancyhead[L]{\small   \textit{Banaras Hindu University}}
\fancyhead[R]{\small   \textit{BPT-602: Solid State Physics}}
\fancyfoot[C]{\small Page  \thepage\ of \pageref{LastPage}}


\newlist{parts}{enumerate}{2}
\setlist[parts,1]{label=(\alph*),leftmargin=*,itemsep=5pt,topsep=3pt}
\setlist[parts,2]{label=(\roman*),leftmargin=*,itemsep=3pt,topsep=2pt}

\newcommand{\pts}[1]{\hfill   \textit{[#1]}}


\begin{document}


\begin{center}
  {\large\bfseries BANARAS HINDU UNIVERSITY}\\[2pt]
  {\normalsize B.Sc. (Hons.) Semester VI Examination, 2022-23}\\[6pt]
  \rule{0.8\linewidth}{0.5pt}\\[6pt]
  {\large\bfseries Paper No. BPT-602: Solid State Physics}\\[4pt]
  {\normalsize Time: 3 Hours\hfill Maximum Marks: 70}
\end{center}
\rule{\linewidth}{0.4pt}
\smallskip



\noindent   \textit{Instructions: Answer all questions. Symbols have their usual meanings.}
\bigskip




\noindent   \textbf{Question 1.} Attempt any   \textbf{four} of the following: \pts{4 $ \times$ 3.5 = 14}

\begin{parts}
  \item What are the main differences in terms of structure among single-crystalline, polycrystalline, quasi-crystalline, and amorphous materials?
  \item State Bragg's law and Laue's conditions for X-ray diffraction. Show that they are equivalent.
  \item What are the differences in the assumptions of the Drude model and the Sommerfeld model for the free electron theory of metals? What do you understand by the 'independent electron approximation'?
  \item What is the difference between mirror and inversion symmetries? Explain with a suitable diagram.
  \item Draw the $E-k$ diagram of a nearly free electron in a 1D periodic potential. How is it different from the $E-k$ diagram of a free electron? Explain briefly.
\end{parts}

\medskip\rule{\linewidth}{0.2pt}




\noindent   \textbf{Question 2.}

\begin{parts}
  \item What do you understand by reciprocal lattice? Establish the relationship between the real lattice and the reciprocal lattice. Describe the Ewald construction. \pts{7}
  \item Show that the reciprocal lattice for an FCC lattice is a BCC lattice. \pts{7}
\end{parts}

\medskip\rule{\linewidth}{0.2pt}




\noindent   \textbf{Question 3.}

\begin{parts}
  \item Draw an FCC lattice and write down the lattice constants of the primitive unit cell. \pts{3.5}
  \item Determine the packing fraction of the FCC lattice. \pts{3.5}
  \item Which is the close-packed plane (CPP) and close-packed direction in the FCC lattice? \pts{3.5}
  \item How many CPPs are there in an FCC lattice? \pts{3.5}
\end{parts}
\hfill   \textbf{OR}

\begin{parts}
  \item Using Sommerfeld theory, find the components of the drift velocity of electrons in the presence of electric and magnetic fields. \pts{7}
  \item What is the Hall field? Derive the relation for the Hall coefficient of a metal. Discuss the utilities of the Hall effect in determining the properties of a material. \pts{7}
\end{parts}

\medskip\rule{\linewidth}{0.2pt}




\noindent   \textbf{Question 4.}

\begin{parts}
  \item Derive the transcendental equation in the Kronig-Penney model. \pts{9}
  \item Which parameters in the Kronig-Penney model determine the magnitude of the band gap? Explain briefly. \pts{5}
\end{parts}
\hfill   \textbf{OR}

\begin{parts}
  \item Derive the dispersion relations for 1-dimensional diatomic lattice vibrations and draw the dispersion curves. Show the atomic vibrations in the optical and acoustical branches at $K  o 0$ and at $K = \pi/2a$ (where '$2a$' is the periodicity). \pts{14}
\end{parts}

\medskip\rule{\linewidth}{0.2pt}




\noindent   \textbf{Question 5.}

\begin{parts}
  \item Define magnon. Establish the dispersion relation for the magnon of a ferromagnet. Show that each spin precesses circularly about the $z$-axis. In the long wavelength region, compare the dispersion relation of the magnon with that of the phonon. Draw these plots. \pts{14}
\end{parts}

\vfill
\rule{\linewidth}{0.4pt}

\begin{center}\small
  \href{https://bhu-exam.vercel.app/}{Click here to visit our website}\\[4pt]
  \href{https://me-aryan.vercel.app/}{Click here to visit our contributor}\\[4pt]
  \href{https://quashan-me.vercel.app/}{Click here to visit our contributor}
\end{center}

\end{document}
